\answer{ex:05:1}{$\ell\approx17\um$؛ و$\approx100$ يوم: البثّ عبر الدماغ يؤديه تفرّع السلك، أما الانتشار فلا يفعل إلا أن يبسط كل موقع تحرير على $\sim20\um$.}
\answer{ex:05:2}{$r=ab/(1+G)$، و$S=1/(1+G)$ تمامًا؛ $+1.7\,\%$ بدل $+20\,\%$.}
\answer{ex:05:3}{$T^*=1.21\,\text{s}$ عند $G=2$ و$0.19\,\text{s}$ عند $G=9$: عمليًا $G\sim2\text{--}4$، أي كبت بـ$3\text{--}5$ مرات.}
\answer{ex:05:4}{$\mathrm{CV}=1/\sqrt{2\lambda\tau_c}\approx9\cdot10^{-4}$ عند تردد كلي $\lambda$؛ والعصبون الواحد يحتاج إلى $\tau_c=12.5\,\text{s}$.}
\answer{ex:05:5}{$c\approx0.40$ في الحالتين؛ $\mathrm{CV}=0.050$ مقابل $0.158$، أي أقل بـ$\sqrt{10}$ مرة. تبقى ترددات العصبونات المفردة متناسبة مع $a_i$: لا يُضبط إلا المجموع.}
\answer{ex:05:6}{$N_1=\overline{A\vee B}$، $N_2=\overline{A\vee N_1}$، $N_3=\overline{B\vee N_1}$، $N_4=\overline{N_2\vee N_3}$، و$\mathrm{XOR}=N_5=\overline{N_4}$؛ كل تبديل يستغرق $\tau_c\ln2$، والمسار من $4$ بوابات يستغرق $2.8\,\text{s}$ لكل \foreignlanguage{english}{XOR}.}
\answer{ex:05:7}{(أ)~$\tau_\gamma=6/\ln3\approx5.5\,\text{s}$. (ب)~$\bar D<4\,\text{s}$ مقابل $D<2.2\,\text{s}$: عدم إمكان التنبؤ بالتأخير يزيد الصبر.}
\answer{ex:05:8}{$k_2=1/\tau_d=10\,\text{s}^{-1}$، $k_1=1/\tau_d-1/\tau_\gamma=9.5\,\text{s}^{-1}$؛ $\tau_\gamma=1/(k_2-mk_1)$: يتضاعف الأفق عند $+2.6\,\%$ (وعند $+5.3\,\%$ يصير لانهائيًا).}
